\subsection{The sequencer}
\label{sec:ha:disabled:sequencer}

The sequencer dies and is restarted. No member's connection is broken by this:
their sessions are held by gateways, which stay up. What stops is the venue's
ability to take an order.

\subsubsection{What state it holds}
\label{sec:ha:disabled:sequencer:state}

\textbf{The ordering.} The number it will assign next, and the write-ahead log
that records what it has assigned. This is the venue's account of what it has
taken, and everything else in this subsection is subordinate to it.

\textbf{The leadership generation.} Which generation of leadership its work is
stamped with, so that a message from a superseded instance can be told from a
current one.

\textbf{Each session's numbering.} Where every session has reached in each
direction, and which of its outbound numbers carried execution reports. The
sequencer holds this on the gateways' behalf and not for itself ---
\cref{sec:ha:disabled:gateway:durable} --- because it is the only component that
every gateway instance talks to.

\textbf{Where an execution report should go.} The route from a committed order
back to the session that placed it.

\textbf{What it has deferred.} How many orders it has taken that no matching
engine has yet seen, and how long that has been true.

\subsubsection{Where the copy that outlives it lives}
\label{sec:ha:disabled:sequencer:durable}

\textbf{The first two survive it.} The log is on disk, and the next number is
taken from its last record. The generation is in a small file beside the log,
written to a temporary file and renamed over the real one, so that a reader sees
either the old value or the new one whatever moment the process died at.

\textbf{The other three have no durable copy.} They are held only in the running
process, and what happens to each on a restart is different.

Session numbering is \emph{relearned} rather than recovered. Gateways report
every session's position on a timer, so a restarted sequencer is told again,
within one reporting interval, about every session that is bound to a gateway at
that moment. A session with no live connection has nobody to tell it --- and
that is precisely the session that is about to reconnect and ask.

\begin{gap}{}
  A session that was not connected when the sequencer restarted is forgotten
  entirely. Nothing on disk holds its numbering and no gateway is reporting it.
\end{gap}

The route from an order to a session is deliberately not rebuilt. Recovering it
by replaying the log would populate it with entries that nothing ever erases,
and the log holds a trading day.

\begin{gap}{}
  An execution report arriving for an order committed before the restart cannot
  be placed, because the route it needed was not kept and is not rebuilt.
\end{gap}

What has been deferred is a count and a time. The orders themselves are in the
log, and a restarted sequencer has no record that any of them is still to be answered.

\begin{gap}{BUG-0064}
  A restart resets the venue's knowledge that it is holding orders no matching
  engine has seen. The records remain in the log; what is lost is that they were
  waiting.
\end{gap}

\subsubsection{How a restarted instance becomes current}
\label{sec:ha:disabled:sequencer:recovery}

It reopens the log, takes the next number from the last record it finds, and
reads the generation from the file beside it. What it cannot recover, it is
told: gateways report the sessions they hold, and a matching engine connects.

What is outstanding at the moment it dies falls into cases, and they are not
alike.

\textbf{An order committed and not yet forwarded to the matching engine} is the
venue's, and is covered by \reqref{R-0013}: the engine reports how far it has
got and the sequencer sends what follows.

\textbf{An execution report produced but not yet sent onward} is a different
problem. The report itself survives, because it was written to the log with the
session it belongs to stamped on it. What does not survive is the route the
sequencer was going to use.

\begin{requirement}{R-0025}{A report can be addressed to its session after a restart}
  An execution report the venue has produced shall be deliverable to the session
  that placed the order, after a restart of any component between the two, and
  the session it belongs to shall be determinable from what survived that
  restart.
  \rationale{The record of the report already names its session. A route held
  only in a running process is lost at precisely the moment it is wanted, and
  nothing else can supply it.}
  \uncovered{no scenario restarts a sequencer with reports undelivered}
\end{requirement}

\textbf{This is not the same operation as a promotion}, and the difference
matters more than it first appears. A follower has been receiving each session
report all along, because gateways report to both instances, and has been
receiving every log record as it was written. It therefore already holds what a
restarted instance must be told --- including for sessions that are not bound to
any gateway, which are exactly the ones nobody will report.

So the sequencer is the reverse of \cref{sec:ha:disabled:gateway}, where there
is only one operation because there is no promotion. Here there are two, they
recover different amounts, and the weaker of them is the one the venue falls
back on when there is nothing to fail over to.

\subsubsection{What the member is told}
\label{sec:ha:disabled:sequencer:member}

\begin{requirement}{R-0026}{A logon completes or is refused within a bounded time}
  A member logging on shall be given either an established session or a refusal,
  within a stated time, whatever the venue has restarted meanwhile. It shall not
  be left authenticated and unestablished.
  \rationale{A session that is neither established nor refused leaves the member
  unable to trade and with nothing to act on: it cannot proceed, and it has been
  given no reason to reconnect.}
  \uncovered{no scenario restarts a sequencer during a logon}
\end{requirement}

\begin{requirement}{R-0027}{An interrupted resend leaves the member able to act}
  Where the venue cannot finish supplying a range of messages it has begun, it
  shall say so. A member shall not be left with an unfilled gap and no
  indication that nothing further is coming.
  \rationale{From the member's side, waiting for records that will never arrive
  is indistinguishable from waiting for records that are merely slow. Only the
  venue can tell the two apart, so only the venue can end it.}
  \uncovered{no scenario restarts a sequencer during a resend}
\end{requirement}

\begin{requirement}{R-0011}{A session's numbering survives a restart of whatever remembers it}
  A member reconnecting after the venue has restarted the component that holds
  its numbering shall be resumed where it left off, and not at the beginning.
  \rationale{A member whose numbering restarts sees a break it cannot
  reconcile, and the venue cannot tell it which of its orders the venue still
  holds.}
  \uncovered{no scenario restarts a sequencer with a session unbound}
\end{requirement}

\begin{requirement}{R-0012}{A member is never resumed at a numbering the venue cannot vouch for}
  Where the venue has no record of a session's numbering, it shall say so rather
  than resume that session at the beginning.
  \rationale{Resuming at the beginning is indistinguishable, to a member whose
  own store has also restarted, from a session that has genuinely never traded.
  Both sides then agree on a numbering neither has any basis for, no gap is
  visible to either, and the member is silently resynchronised while its orders
  are still open. Being refused is recoverable; being agreed with wrongly
  is not.}
  \uncovered{no scenario restarts a sequencer with a session unbound}
\end{requirement}

\begin{gap}{}
  A member whose session the venue has forgotten is resumed at the beginning and
  is not told.
\end{gap}

\begin{requirement}{R-0013}{An order committed before a sequencer restart still receives its execution report}
  An order that the sequencer had committed before it died shall be processed,
  and an execution report shall be produced for it and be deliverable to the
  session that placed the order.
  \rationale{Commitment is the venue taking the order, and a restart of the
  component that took it does not give it back.}
  \uncovered{no scenario restarts a sequencer with orders committed and unanswered}
\end{requirement}

\begin{requirement}{R-0014}{A sequence number is never issued twice}
  A restarted sequencer shall not assign a number it has assigned before.
  \rationale{The number is the venue's name for an order. Two orders sharing one
  cannot both be answered, and the second cannot be distinguished from a
  retransmission of the first.}
  \uncovered{scenario 14 covers a restart with a peer to synchronise from}
\end{requirement}
